\section{The nine I/O backends}
\label{sec:backends}

Table~\ref{tab:modes} summarises the nine backends. Seven base threads are present in
every mode: main/render, the stdin reader, the logger, the store, the sender, the
listener and metrics.

\begin{table}[h]
\centering\small
\begin{tabularx}{\linewidth}{@{}lcc>{\raggedright\arraybackslash}Xlc@{}}
\toprule
\code{-{}-io=} & conn./peer & fds/peer & read path & write path & threads \\
\midrule
\mode{threads}      & 2 & 2 & blocking reader thread per connection & blocking, on the sender & $8+N$ \\
\mode{pool-sender}  & 2 & 2 & as \mode{threads} & non-blocking pool + pump & $9+N$ \\
\mode{(e)poll} & 2 & 2 & one pool, one I/O thread & same pool & 8 \\
\mode{(e)poll-duplex} & 1 & 1 & one pool, one I/O thread & same pool, same socket & 8 \\
\mode{(e)poll-split} & 1 & 1 & rx pool + I/O thread & \code{TxPoller} + tx thread & 9 \\
\mode{peer-threads} & 1 & 2 & one owner thread per peer & one-connection pool & $8+N$ \\
\bottomrule
\end{tabularx}
\caption{Backends. \mode{(e)poll} stands for both the \code{poll} and \code{epoll} variant. $N$ is the number of peers. ``fds/peer'' excludes the node's fixed descriptors.}
\label{tab:modes}
\end{table}

\subsection{\mode{threads} and \mode{pool-sender}}

\code{ListenerWorker} accepts connections in a blocking loop and starts one
\code{ReaderWorker} thread per connection. Each reader waits in \code{receive()} with a
200\,ms timeout and bumps a heartbeat counter.

\code{WatchdogWorker} looks after the readers:
\begin{itemize}
  \item It scans every 100\,ms and joins readers that have finished.
  \item A reader whose heartbeat has not moved for 10 scans (1\,s) is treated as
        wedged; the watchdog shuts its socket down to unblock it.
  \item It enforces \code{-{}-max-peers}. Once the cap is exceeded, the least recently
        active readers are evicted first.
\end{itemize}

The two modes differ only in how they write:
\begin{itemize}
  \item \mode{threads} writes with blocking \code{send}s on the sender thread. A peer
        that stops reading stalls that thread for up to 500\,ms per attempt.
  \item \mode{pool-sender} still dials with a blocking \code{connect}, but then hands
        the socket to a \code{TcpConnectionPool}. Every send is non-blocking, and
        whatever doesn't fit is queued, up to 1\,MiB. A pump thread drives the pool,
        flushing queued bytes and noticing dead sockets.
\end{itemize}

\subsection{\mode{poll} and \mode{epoll}, simplex and duplex}

All four modes are \code{MultiplexedIo<Pool>}: one pool holds every connection, and
one I/O thread loops on \code{pool.receive()}. A single call returns bytes from any
readable connection and flushes queued writes on any writable one. \code{POLLOUT}
(or \code{EPOLLOUT}) is armed only while a transmit buffer holds data.

\code{registry\_} maps pool ids to \code{FrameChannel}s. Whoever adds a connection
holds \code{registry\_mutex\_} across both the pool insert and the registry insert, so
the I/O thread can never receive bytes for an id it cannot find.

The two pools differ in how they track sockets:
\begin{itemize}
  \item \code{TcpConnectionPool} keeps a desired-state table and rebuilds its
        \code{pollfd} array before every \code{poll()}. Other threads wake it through
        an \code{eventfd(2)}. A round-robin cursor decides which ready socket is
        served first.
  \item \code{EpollConnectionPool} calls \code{epoll\_ctl} directly, even while
        another thread is blocked in \code{epoll\_wait}. Its wake-up eventfd uses the
        reserved token 0 (\code{kWakeToken}), and it keeps a fairness cursor over the
        ready events.
\end{itemize}

The \emph{Dual} model (\mode{poll}, \mode{epoll}) keeps two connections per peer in
one pool, with capacity $2\times$\code{max\_peers}. Accepted sockets have a channel;
dialled ones carry writes only. The \emph{Duplex} model gives every connection a
channel, whether it was accepted or dialled, and the \code{ConnectAck} goes back on
the same handle. Its capacity is \code{max\_peers}${}+16$. Duplex halves the sockets
and avoids the dial-back.

\subsection{\mode{poll-split} and \mode{epoll-split}}

These modes use the same duplex link logic, but read and write each socket on
different threads:
\begin{itemize}
  \item Each connection is a \code{shared\_ptr<TcpConnection>} held by three owners:
        the rx pool (which only reads it), a \code{TxPoller} (which only writes it)
        and the registry.
  \item The \code{TxPoller} has its own thread and waits for writability only on
        connections that have queued bytes.
  \item Socket options are applied eagerly (\code{apply\_options\_now()}) before the
        connection is shared, so neither thread changes them later.
  \item A fatal write calls \code{shutdown(SHUT\_RDWR)}. The rx pool then sees EOF and
        reports the disconnect, so disconnects are always reported from the read
        side.
\end{itemize}

\subsection{\mode{peer-threads}}

Each peer is a \emph{link}: a one-connection \code{TcpConnectionPool} driven by its
own owner thread (\code{pt-<handle>}).
\begin{itemize}
  \item \textbf{Reads and flushing:} the owner thread reads the socket and flushes
        queued writes.
  \item \textbf{Immediate writes:} the shared sender writes through
        \code{pool.send()} under the pool's mutex, so frames never interleave.
  \item \textbf{Ending a link:} the owner thread is the only caller of
        \code{on\_closed()}, so the link is reported gone exactly once. That holds
        whichever of the four triggers comes first: the peer closing, a malformed
        frame, the sender's \code{close()}, or a watchdog eviction.
  \item \textbf{Cap:} links count against \code{-{}-max-peers} through the shared
        watchdog, which evicts the link with the oldest \code{last\_read\_ns}. Stall
        detection is turned off, because an idle owner legitimately blocks in
        \code{receive()}.
  \item \textbf{No restarts:} an evicted owner thread is joined and gone. The peer
        can only come back as a new link, with a new handle and \code{connection\_id}.
\end{itemize}

\textbf{Limitations.} Retries are a 100\,ms sleep on the shared sender thread, and
\code{open()} performs a blocking \code{connect()} on that same thread.
Section~\ref{sec:res-churn} measures the cost of the first.
